\section*{الفصل \ref{ch:g12:nuclear-energy} --- الطاقة النووية: الانشطار والاندماج و$E=mc^2$}

\begin{solution}{exo:g12:nuclear-energy:1}
(أ) $\qty{1}{eV} = \qty{1.60e-19}{J}$؛
$\qty{1}{MeV} = \qty{1.60e-13}{J}$.
(ب) $200 \times \num{1.60e-13} = \qty{3.2e-11}{J}$.
(ج) $E = \num{1.66054e-27} \times \num{9.00e16} = \qty{1.49e-10}{J}
= \num{1.49e-10}/\num{1.60e-13} \approx \qty{9.3e2}{MeV}$ --- أي
\qty{931.5}{MeV} المعتمدة في الدرس، في حدود تدويرنا للمقدار $c$.
\end{solution}

\begin{solution}{exo:g12:nuclear-energy:2}
$E = \num{6.0e-3} \times \num{9.00e16} = \qty{5.4e14}{J}$؛
$\num{5.4e14}/\num{3.6e7} = \num{1.5e7}$ يوم $\approx \num{4.1e4}$
سنة --- أي أربعون ألفية على مكعب سكر واحد.
\end{solution}

\begin{solution}{exo:g12:nuclear-energy:3}
$\Delta m = 6 \times 1.00728 + 6 \times 1.00866 - 11.99671 =
\qty{0.09893}{u}$؛ $E_b = 0.09893 \times 931.5 \approx \qty{92.2}{MeV}$؛
$E_b/A \approx \qty{7.68}{MeV}$ لكل \omterm{def:g11:fundamental-interactions:ladder}{نوكليون}.
\end{solution}

\begin{solution}{exo:g12:nuclear-energy:4}
(أ) $A = 236 - 144 - 3 = 89$، $Z = 92 - 56 = 36$:
${}^{89}_{36}\mathrm{Kr}$، أي الكريبتون.
(ب) $A = 236 - 97 - 2 = 137$، $Z = 92 - 40 = 52$:
${}^{137}_{52}\mathrm{Te}$، أي التيلوريوم.
\end{solution}

\begin{solution}{exo:g12:nuclear-energy:5}
نحو $7.1$ و$7.7$ و$8.8$ و\qty{7.6}{MeV} لكل \omterm{def:g11:fundamental-interactions:ladder}{نوكليون}؛ والحديد-56 هو
الأشدّ ارتباطًا. واليورانيوم يقدر أن يحرّر \omterm{def:g10:energy-conservation:energy}{طاقة} \omterm{def:g12:nuclear-energy:fission}{بالانشطار}، والهيليوم-4 و
الكربون-12 \omterm{def:g12:nuclear-energy:fusion}{بالاندماج} --- فكلتا الحركتين تصعد نحو الحديد؛ أما الحديد نفسه فليس
له مذهب.
\end{solution}

\begin{solution}{exo:g12:nuclear-energy:6}
$\Delta m = 1.00728 + 1.00866 - 2.01355 = \qty{0.00239}{u}$، ومنه
$E_b \approx \qty{2.23}{MeV}$ --- أي نحو مليون ضعف من بضعة
\unit{eV} في رابطة كيميائية.
\end{solution}

\begin{solution}{exo:g12:nuclear-energy:7}
قبل: $\qty{236.00216}{u}$؛ وبعد: $143.89220 + 88.89790 +
3 \times 1.00866 = \qty{235.81608}{u}$؛ $\Delta m = \qty{0.18608}{u}$،
ومنه $E \approx \qty{173}{MeV} = 173 \times \num{1.60e-13} \approx
\qty{2.8e-11}{J}$.
\end{solution}

\begin{solution}{exo:g12:nuclear-energy:8}
$\Delta m = 2 \times 2.01355 - 3.01493 - 1.00866 = \qty{0.00351}{u}$:
$E \approx \qty{3.27}{MeV}$. \omterm{def:g12:nuclear-energy:fusion}{فاندماج} D--T يصنع هيليوم-4، وهو مرتبط ارتباطًا شاذًّا
(\qty{7.1}{MeV} لكل \omterm{def:g11:fundamental-interactions:ladder}{نوكليون})؛ أما الهيليوم-3 فيجلس أخفض بكثير على المنحني، ومن ثم
تحرّر الخطوة أقلّ بكثير.
\end{solution}

\begin{solution}{exo:g12:nuclear-energy:9}
(أ) النيوترون متعادل: فلا حاجز كولوم، \omterm{def:g12:nuclear-energy:fission}{والانشطار} يشتعل على البارد.
أما النواتان المندمجتان فموجبتان كلتاهما وعليهما تسلّق الحاجز، ومن ثم لا يعطي
تصادمات عنيفة بما يكفي إلا اضطراب قدره $\sim \qty{1e8}{K}$. (ب) و
قلب الشمس يحوي عددًا فلكيًّا من البروتونات؛ وأسرعها في الحشد، وهي نادرة،
هي التي تندمج، وأمام الشمس مليارات السنين تنفقها
--- فالاحتراق البطيء هو بالضبط ما يريده نجم. (ج) ولا وجود لانهيار
نيوتروني: فكل \omterm{def:g12:nuclear-energy:fusion}{اندماج} يشتريه تصادمه الخاص، وأي فقدان
للحبس يبرّد \omterm{def:g12:nuclear-energy:barrier}{البلازما} ويطفئ النار.
\end{solution}

\begin{solution}{exo:g12:nuclear-energy:10}
$\Delta m = 4 \times 1.00728 - 4.00151 = \qty{0.02761}{u}$، ومنه
$E \approx \qty{25.7}{MeV}$؛ والجزء
$0.02761/4.02912 \approx 0.7\%$ من الكتلة.
\end{solution}

\begin{solution}{exo:g12:nuclear-energy:11}
(أ) يبطّئ نيوترونات \omterm{def:g12:nuclear-energy:fission}{الانشطار} بالتصادمات؛ والنيوترونات البطيئة أرجح
كثيرًا في شطر اليورانيوم-235، ومن ثم تحتاجها السلسلة. (ب) وهي
تمتصّ النيوترونات، فتمسك السلسلة عند نيوترون واحد لكل \omterm{def:g12:nuclear-energy:fission}{انشطار} ---
فإذا أُدخلت ماتت. (ج) ولا ماء، فلا تهدئة: إذ تبقى النيوترونات
سريعة، فتُخطئ، وتختنق السلسلة --- فالتصميم يفشل نحو
``الإطفاء''.
\end{solution}

\begin{solution}{exo:g12:nuclear-energy:12}
تتضاعف الانشطارات في كل جيل، ومن ثم يكون قد انشطر بعد $n$ جيلًا نحو $2^n$؛
ويعطي $2^{n} = \num{2.56e21} \approx 2 \times (10^{3})^{7}
\approx 2^{71}$ المقدارَ $n \approx 71$. والزمن: $71 \times \qty{10}{ns}
\approx \qty{0.7}{\micro s}$ --- أي غرام من الوقود، بما يساوي ثلاثة أطنان من الفحم،
في أقل من ميكروثانية: فالحرجية هاوية، لا منحدر.
\end{solution}

\begin{solution}{exo:g12:nuclear-energy:13}
$\Delta m = 55.92066 - 2 \times 27.96924 = \qty{-0.01782}{u}$:
$E \approx \qty{-16.6}{MeV}$ --- أي إن الشطر \emph{يمتصّ} \omterm{def:g10:energy-conservation:energy}{طاقة}.
فالحديد يجلس عند الذروة: وشطره أو دمجه كلاهما ينزل في
الارتباط، وكلاهما يكلّف \omterm{def:g10:energy-conservation:energy}{طاقة}. وهو رماد كل نار نووية.
\end{solution}

\begin{solution}{exo:g12:nuclear-energy:14}
الحديد: $\Delta m = 26 \times 1.00728 + 30 \times 1.00866 - 55.92066 =
\qty{0.52842}{u}$، $E_b \approx \qty{492}{MeV}$، $E_b/A \approx
\qty{8.79}{MeV}$. واليورانيوم: $\Delta m = 92 \times 1.00728 + 143 \times
1.00866 - 234.99350 = \qty{1.91464}{u}$، $E_b \approx \qty{1784}{MeV}$،
$E_b/A \approx \qty{7.59}{MeV}$. وإعادة رصّ $235$ \omterm{def:g11:fundamental-interactions:ladder}{نوكليون} عند
\qty{8.5}{MeV} بدلًا من \qty{7.59}{MeV} تكسب $235 \times 0.9
\approx \qty{2e2}{MeV}$ --- أي $185$--$200$ في المثال المحلول.
\end{solution}

\begin{solution}{exo:g12:nuclear-energy:15}
$\num{5.0e8}/\num{2.82e-12} \approx \num{1.8e20}$ تفاعل في
\omterm{def:g10:orders-of-magnitude:unit}{الثانية}؛ ومعدّل الكتلة $\num{1.8e20} \times 5.029 \times \num{1.66054e-27}
\approx \qty{1.5e-6}{kg/s}$، أي\ نحو \qty{0.13}{kg} في اليوم ---
في مقابل \num{1.4e6} \omterm{def:g10:orders-of-magnitude:unit}{كيلوغرام} من الفحم: أي عشرة ملايين إلى واحد.
\end{solution}

\begin{solution}{pb:g12:nuclear-energy:1}
\textbf{1.} $P_{\text{حر}} = \num{1.0e9}/0.33 \approx \qty{3.0}{GW}$.

\textbf{2.} $200 \times \num{1.60e-13} = \qty{3.2e-11}{J}$.

\textbf{3.} $\num{3.0e9}/\num{3.2e-11} \approx \num{9.4e19}$ \omterm{def:g12:nuclear-energy:fission}{انشطار}
في \omterm{def:g10:orders-of-magnitude:unit}{الثانية}.

\textbf{4.} $m_{\mathrm U} = 235 \times \num{1.66054e-27} =
\qty{3.90e-25}{kg}$؛ $\num{9.4e19} \times \num{3.90e-25} \approx
\qty{3.7e-5}{kg/s}$ --- أي نحو أربعين ميكروغرامًا في \omterm{def:g10:orders-of-magnitude:unit}{الثانية}.

\textbf{5.} $\num{3.7e-5} \times \num{3.16e7} \approx \qty{1.2e3}{kg}$:
أي نحو $1.2$ طن من اليورانيوم-235 في السنة.

\textbf{6.} $\num{1.2e3}/0.040 \approx \qty{2.9e4}{kg}$ --- أي نحو $29$
طن من الوقود: أي شاحنة ثقيلة واحدة، مرة في السنة.

\textbf{7.} $\num{3.0e9} \times \num{3.16e7} \approx \qty{9.5e16}{J}$.

\textbf{8.} $\num{9.5e16}/\num{3.0e7} \approx \qty{3.2e9}{kg}$: أي $3.2$
مليون طن من الفحم.

\textbf{9.} $\num{3.2e9}/\num{1.2e3} \approx \num{2.7e6}$ --- أي
$\num{8.2e13}/\num{3.0e7}$ في السُّلّم، كما يجب.

\textbf{10.} $\num{3.2e9}/\num{3.0e6} \approx 1060$ قطارًا في السنة،
أي نحو ثلاثة في اليوم: فقطار فحم عند الفجر والظهر والغسق إلى الأبد،
في مقابل شاحنة واحدة كل ربيع.

\textbf{11.} $17.6 \times \num{1.60e-13} = \qty{2.82e-12}{J}$.

\textbf{12.} الزوج الواحد يزن $5.029 \times \num{1.66054e-27} =
\qty{8.35e-27}{kg}$: $\num{2.82e-12}/\num{8.35e-27} \approx
\qty{3.4e14}{J/kg}$ --- أي درجة \omterm{def:g12:nuclear-energy:fusion}{الاندماج} في السُّلّم.

\textbf{13.} $\num{9.5e16}/\num{3.4e14} \approx \qty{2.8e2}{kg}$ من
وقود D--T؛ وحصة الديوتيريوم $280 \times 2.014/5.029 \approx \qty{1.1e2}{kg}$.

\textbf{14.} $\num{1.1e2}/\num{3.3e-2} \approx \qty{3.4e3}{m^3}$ من
ماء البحر --- أي نحو مسبح ونصف لسنة المدينة كلها.

\textbf{15.} لا شيء يتضاعف: فلا نيوترون يشعل \omterm{def:g12:nuclear-energy:fusion}{الاندماج} التالي،
وفقدان الحبس يبرّد \omterm{def:g12:nuclear-energy:barrier}{البلازما} ويوقفها في ثوانٍ.
والجزء الصعب هو النقيض --- أي إمساك \qty{1.5e8}{K} معًا مدةً
تكفي للاحتراق.

\textbf{16.} $\num{3.9e26}/\num{9.00e16} \approx \qty{4.3e9}{kg/s}$:
أي أربعة ملايين طن في \omterm{def:g10:orders-of-magnitude:unit}{الثانية}.

\textbf{17.} $\num{3.0e9}/\num{9.00e16} \approx \qty{3.3e-8}{kg/s}$،
أي\ $\num{3.3e-8} \times \num{3.16e7} \approx \qty{1.0}{kg}$ في السنة
--- ونعم: $\qty{1.2e3}{kg} \times 0.00085 \approx 1$، $\num{3.2e9}
\times \num{3.3e-10} \approx 1$، $280 \times 0.00375 \approx 1$. فكل
وقود يسلّم \omterm{def:g10:orders-of-magnitude:unit}{الكيلوغرام} نفسه؛ ولا يختلفون إلا في مقدار الحمولة
التي يجب أن تحمله.

\textbf{18.} $\num{4.3e9} \times \num{3.16e7} \times \num{4.5e9}
\approx \qty{6.1e26}{kg}$ --- أي $\num{6.1e26}/\num{2.0e30} \approx
\num{3e-4}$ لا غير، أي ثلاثة أجزاء في عشرة آلاف.

\textbf{19.} الكتلة القابلة للتحويل $0.0007 \times \num{2.0e30} =
\qty{1.4e27}{kg}$، وتساوي $\num{1.4e27} \times \num{9.00e16} \approx
\qty{1.3e44}{J}$؛ وعند \qty{3.9e26}{W} تدوم $\num{1.3e44}/
\num{3.9e26} \approx \qty{3.2e17}{s} \approx \num{1.0e10}$ سنة ---
أي عشرة مليارات، منها نحو خمسة مليارات باقية.

\textbf{20.} سنة واحدة من المدينة: $3.2$ مليون طن من الفحم، أو
$\num{1.2e3}$ \omterm{def:g10:orders-of-magnitude:unit}{كيلوغرام} من اليورانيوم-235، أو $\num{2.8e2}$ \omterm{def:g10:orders-of-magnitude:unit}{كيلوغرام} من
وقود D--T مسحوبًا من مسبح ونصف من ماء البحر. وتحت كل دفتر
القيد نفسه: \omterm{def:g10:orders-of-magnitude:unit}{كيلوغرام} واحد من الكتلة، ذهب. و$E = mc^2$ لا يعبأ بأي
وقود يحمله --- بل بحجم الشاحنة التي يلزمه.
\end{solution}
